%% kapa: exam=1 round=27 year=2016 period=1 form=A subject=real_estate subject_ko=부동산학원론 seq=16 number=96
%% source: Q-Net 감정평가사 기출문제 (cst00309) artlSeq=5169402 file=extracted/1008/5169402/제27회 감정평가사 1차 문제/제27회 감정평가사 1교시문제(A형).pdf
%% conversion: kordoc 4.24.51 -> tools/gen1.py (auto); review=proofread-sonnet; flags=none
\begin{question}{96}
도시 및 부동산개발에 관한 설명으로 옳지 않은 것은?
\begin{choices}
  \item 부동산개발업의 관리 및 육성에 관한 법령상 부동산개발이란 토지를 건설공사의 수행 또는 형질변경의 방법으로 조성하면서 시공을 담당하는 행위를 말한다.
  \item 부동산개발업의 관리 및 육성에 관한 법령상 부동산개발업이란 타인에게 공급할 목적으로 부동산개발을 수행하는 업을 말한다.
  \item 부동산개발업의 관리 및 육성에 관한 법령상 공급이란 부동산개발을 수행하여 그 행위로 조성ㆍ건축ㆍ대수선ㆍ리모델링ㆍ용도변경 또는 설치되거나 될 예정인 부동산, 그 부동산의 이용권으로서 대통령령으로 정하는 권리의 전부 또는 일부를 타인에게 판매 또는 임대하는 행위를 말한다.
  \item 도시개발법령상 도시개발사업이란 도시개발구역에서 주거, 상업, 산업, 유통, 정보통신, 생태, 문화, 보건 및 복지 등의 기능이 있는 단지 또는 시가지를 조성하기 위하여 시행하는 사업을 말한다.
  \item 도시 및 주거환경정비법령상 건축물이 훼손되거나 일부가 멸실되어 붕괴 그 밖의 안전사고의 우려가 있는 건축물은 노후ㆍ불량건축물에 해당한다.
\end{choices}
\end{question}
\answer{1}
